\subsection*{導入：なぜソフトウェアアーキテクチャを独立して学ぶのか}
\addcontentsline{toc}{subsection}{導入：なぜソフトウェアアーキテクチャを独立して学ぶのか}
SWEBOK v4.0a では、ソフトウェアアーキテクチャをソフトウェア設計から独立した知識領域として扱う。これは、1990 年代以降にアーキテクチャ研究と実務が大きく発展し、設計の一部としてだけでは説明しきれない知識が増えたためである。

ソフトウェアアーキテクチャは、概念、表現、作業成果物、文脈、過程、方法、分析、評価という複数の視点から理解する必要がある。特に大規模なシステムでは、個々の実装詳細よりも先に、主要な構成要素、構成要素間の関係、外部環境との相互作用、品質特性への対応、変更に対する耐性を考える必要がある。

\begin{statementbox}{重要}
アーキテクチャは「実装の前に一度だけ描く図」ではない。要求、品質、組織、運用、保守、将来の進化に関わる継続的な判断の集合である。
\end{statementbox}
